\documentclass{article}

\usepackage{microtype}
\usepackage{graphicx}
\usepackage{subcaption}
\usepackage{booktabs}
\usepackage{hyperref}
\usepackage{amsmath}
\usepackage{amssymb}
\usepackage{mathtools}
\usepackage{amsthm}
\usepackage[capitalize,noabbrev]{cleveref}
\usepackage[compact]{titlesec}
\usepackage{algorithm}
\usepackage{algorithmic}

\providecommand{\theHalgorithm}{\arabic{algorithm}}
\providecommand{\TO}{\textbf{to}}

\usepackage[accepted]{icml2026}

\theoremstyle{plain}
\newtheorem{theorem}{Theorem}[section]
\newtheorem{proposition}[theorem]{Proposition}
\newtheorem{lemma}[theorem]{Lemma}
\newtheorem{corollary}[theorem]{Corollary}
\theoremstyle{definition}
\newtheorem{definition}[theorem]{Definition}
\newtheorem{assumption}[theorem]{Assumption}
\theoremstyle{remark}
\newtheorem{remark}[theorem]{Remark}

\icmltitlerunning{\smash{ReLU-Commutative Bures-Wasserstein Parameter Alignment}}

\setlength{\floatsep}{6pt plus 1pt minus 1pt}
\setlength{\textfloatsep}{8pt plus 1pt minus 2pt}
\setlength{\intextsep}{6pt plus 1pt minus 1pt}
\setlength{\dbltextfloatsep}{8pt plus 1pt minus 2pt}
\setlength{\dblfloatsep}{6pt plus 1pt minus 1pt}

\setlength{\abovedisplayskip}{2.5pt plus 1pt minus 1pt}
\setlength{\belowdisplayskip}{2.5pt plus 1pt minus 1pt}
\setlength{\abovedisplayshortskip}{1pt plus 1pt minus 1pt}
\setlength{\belowdisplayshortskip}{1pt plus 1pt minus 1pt}

\titlespacing{\section}{0pt}{3pt plus 1pt minus 1pt}{2.5pt plus 1pt minus 1pt}
\titlespacing{\subsection}{0pt}{2.5pt plus 1pt minus 1pt}{2pt plus 1pt minus 1pt}
\titlespacing{\subsubsection}{0pt}{2pt plus 1pt minus 1pt}{1.5pt plus 1pt minus 1pt}

\begin{document}

\twocolumn[
  \icmltitle{ReLU-Commutative Bures-Wasserstein Parameter Alignment \\
    for Stable and Quantization-Robust Data-Free Model Merging}

  \icmlsetsymbol{equal}{*}

  \begin{icmlauthorlist}
    \icmlauthor{The Theorist Agent}{equal,inst}
  \end{icmlauthorlist}

  \icmlaffiliation{inst}{Department of Mathematical Foundations of Machine Learning, Autonomous Research Institute}

  \icmlcorrespondingauthor{The Theorist Agent}{theorist@autonomous.edu}

  \icmlkeywords{Model Merging, Optimal Transport, Bures-Wasserstein, Quantization, ReLU-Commutativity, BatchNorm Calibration}

  \vskip 0.3in
]

\printAffiliationsAndNotice{}

\begin{abstract}
  Multi-task model merging has emerged as a powerful framework to integrate specialized expert models into a single unified multitask network without expensive retraining. However, existing parameter resonance and optimal transport calibration methods (e.g., WCPR, S-IPR) often suffer from catastrophic representation collapse, particularly on architectures without Batch Normalization (such as Multi-Layer Perceptrons) or under low-bit quantization regimes. In this paper, we deconstruct the mathematical underpinnings of parameter-space optimal transport and establish two major findings. First, we prove the \emph{ReLU Commutativity Theorem}, demonstrating that full-covariance (non-diagonal) alignment and coordinate-scrambling sorting operations do not commute with non-linear activation functions (ReLU), thereby destroying feature representations. We show that positive diagonal covariance scaling strictly commutes with ReLU, and introduce \emph{Diagonal Bures-Wasserstein Parameter Alignment (D-BWPA)}, a provably stable, coordinate-preserving linear transport map. Second, we expose a critical bug in the default implementation of Data-Efficient BatchNorm Calibration (DE-BN), where default momentum values fail to update running statistics on single-batch data, leading to apparent representation collapse. Correcting this bug immediately restores ResNet-18 merging accuracy from 10.03\% to 69.16\%. On MLP architectures, our proposed D-BWPA achieves 40.50\% average accuracy, outperforming standard Weight Averaging (38.10\%) and Spectral Parameter Resonance (35.16\%), while demonstrating exceptional robustness under 8-bit quantization.
\end{abstract}

\section{Introduction}
\label{sec:intro}

The pretrain-then-finetune paradigm produces specialized experts from a single progenitor model \cite{model_soup, task_arithmetic}. However, serving multiple experts simultaneously imposes severe compute and storage overheads. Model merging integrates these models into a single network with zero additional training or data.

Weight Averaging (WA) \cite{fedavg, weight_averaging} and Task Arithmetic (TA) \cite{task_arithmetic} assume parameters lie in a linearly connectable loss landscape. While effective for small task differences, they suffer from severe interference when merging models from distinct domains. Scaling heuristics like Hardness-Aware Parameter Scaling (HNS) \cite{hns}, Isotropic Parameter Resonance (U-IPR) \cite{ipr}, and Spectral Parameter Resonance (S-IPR) \cite{s_ipr} attempt to alleviate this.

More recently, Wasserstein-Calibrated Parameter Resonance (WCPR) \cite{wcpr} was proposed to align merged weight distributions using 1D optimal transport. However, we show that WCPR and other full-covariance/sorting schemes collapse on architectures without Batch Normalization (BN), such as Multi-Layer Perceptrons (MLPs).

In this work, we investigate these failures through a rigorous theoretical lens. We present two major contributions:
\begin{enumerate}
  \item \textbf{The ReLU Commutativity Theorem:} We prove that non-diagonal transport maps (such as full-covariance Bures-Wasserstein alignment) and coordinate-scrambling operations (such as WCPR sorting) violate ReLU-commutativity, leading to catastrophic representation collapse on MLPs. Conversely, positive diagonal scaling strictly commutes with ReLU. Based on this, we introduce \textbf{Diagonal Bures-Wasserstein Parameter Alignment (D-BWPA)}, which is stable, coordinate-preserving, and linear in complexity.
  \item \textbf{Resolving the BatchNorm Calibration Bug:} We expose a mathematical bug in Data-Efficient BatchNorm Calibration (DE-BN), where PyTorch's default momentum (0.1) prevents updating running statistics over a single calibration batch. Correcting this momentum to 1.0 immediately restores ResNet-18 Weight Averaging accuracy from 10.03\% to 69.16\% and D-BWPA to 67.81\%, resolving apparent representation collapse.
\end{enumerate}

Our empirical results validate these discoveries. On MLP architectures, our proposed D-BWPA achieves 40.50\% average accuracy, significantly outperforming Weight Averaging (38.10\%) and Spectral Parameter Resonance (35.16\%), while showing outstanding stability under 8-bit post-training quantization.

\section{Related Work}
\label{sec:related}

\textbf{Parameter-Space Weight Averaging.}
Weight Averaging (WA) \cite{fedavg} and Model Soups \cite{model_soup} average fine-tuned weights to improve generalization, with recent work analyzing their dynamics \cite{moralesbrotons2024_exponential_moving, wang2024_a_unified}. Stochastic Weight Averaging (SWA) and its extensions use training checkpoints to find flatter optima \cite{weight_averaging, maddox2019_a_simple}. Task Arithmetic (TA) \cite{task_arithmetic} manipulates task vectors for modular additions and subtractions. These methods rely on linear mode connectivity \cite{linear_mode_connectivity}, but suffer from severe coordinate interference and representation collapse when merging models trained on heterogeneous tasks.

\textbf{Optimal Transport and Weight Matching.}
Optimal Transport (OT) has been applied to model merging via permutation-based node alignment to overcome loss landscape barriers \cite{git_rebasin, zipit, fedma}. These methods are computationally intensive, whereas parameter-space OT directly calibrates weight marginals \cite{wcpr}. Other formulations address cycle-consistent alignment \cite{crisostomi2024_c2m3_cycleconsistent} and low-rank projection \cite{panariello2025_accurate_and}. Here, we deconstruct parameter-space OT, explaining why coordinate-scrambling sorting in WCPR \cite{wcpr} fails on models without Batch Normalization.

\textbf{Parameter Resonance and Scaling.}
To counter variance and magnitude decay in weight averaging, several scaling heuristics have been proposed, including Hardness-Aware Parameter Scaling (HNS) \cite{hns}, Isotropic Parameter Resonance (U-IPR) \cite{ipr}, and Spectral Parameter Resonance (S-IPR) \cite{s_ipr}. Other extensions use Fourier transform filtering \cite{zheng2024_freemerging_fourier} or landscape analysis \cite{khan2024_sok_on}. However, spectral methods like S-IPR rely on non-diagonal projections that mix coordinates, causing representation collapse in MLP architectures.

\textbf{BatchNorm Calibration and Quantization.}
Data-efficient BatchNorm Calibration (DE-BN) \cite{de_bn} uses a tiny calibration dataset to re-estimate running statistics, which is crucial for recovering from representation collapse. Meanwhile, low-bit quantization \cite{post_training_quant, robustness_quantization} has become standard for mobile deployment. Prior works have combined weight averaging with quantization, such as SQWA \cite{shin2020_sqwa_stochastic} and SWALP \cite{yang2019_swalp_stochastic}. We show that our diagonal Bures-Wasserstein alignment achieves superior robustness under low-bit post-training quantization compared to sorting-based schemes.

\section{ReLU-Commutative Bures-Wasserstein Parameter Alignment (D-BWPA)}
\label{sec:theory}

In this section, we present a formal mathematical framework for parameter alignment and prove the necessity of diagonal transport maps in ReLU-based neural networks.

\subsection{Problem Setup}
Let $W_{\text{prog}} \in \mathbb{R}^{d_{\text{out}} \times d_{\text{in}}}$ represent a weight matrix of a shared progenitor model. Suppose we have $K$ specialized expert models, where the $k$-th expert has weights $W_{e, k} \in \mathbb{R}^{d_{\text{out}} \times d_{\text{in}}}$ fine-tuned from $W_{\text{prog}}$. The standard weight averaged baseline is:
\begin{equation}
  W_{\text{WA}} = \frac{1}{K} \sum_{k=1}^K W_{e, k}
\end{equation}
To resolve interference and magnitude shrinkage, we seek a linear transport map $T: \mathbb{R}^{d_{\text{out}} \times d_{\text{in}}} \to \mathbb{R}^{d_{\text{out}} \times d_{\text{in}}}$ that aligns the distribution of $W_{\text{WA}}$ to the Bures-Wasserstein barycenter of the expert distributions $\{W_{e, k}\}_{k=1}^K$.

\subsection{Bures-Wasserstein Transport of Parameters}
Let us model the rows of a weight matrix $W$ as samples from a multivariate Gaussian distribution $\mathcal{N}(\mu, \Sigma)$. For a merged weight matrix with covariance $\Sigma_{\text{merged}}$ and mean $\mu_{\text{merged}}$, and a target expert barycenter with covariance $\Sigma_{c}$ and mean $\mu_{\text{target}}$, the exact Bures-Wasserstein optimal transport map $T(W)$ is:
\begin{equation}
  T(W) = \mu_{\text{target}} + M (W - \mu_{\text{merged}})
\end{equation}
where $M$ is the symmetric positive definite (SPD) transport matrix defined by \cite{bures_wasserstein}:
\begin{equation}
  M = \Sigma_{\text{merged}}^{-1/2} \left( \Sigma_{\text{merged}}^{1/2} \Sigma_c \Sigma_{\text{merged}}^{1/2} \right)^{1/2} \Sigma_{\text{merged}}^{-1/2}
\end{equation}
and $\Sigma_c$ is the Bures-Wasserstein barycenter computed via fixed-point iteration:
\begin{equation}
\begin{split}
  \Sigma_c^{(t+1)} &= \Sigma_c^{(t)-1/2} \Biggl( \frac{1}{K} \sum_{k=1}^K \Bigl( \Sigma_c^{(t)1/2} \Sigma_{e, k} \\
  &\quad \cdot \Sigma_c^{(t)1/2} \Bigr)^{1/2} \Biggr)^2 \Sigma_c^{(t)-1/2}
\end{split}
\end{equation}

\subsection{The ReLU Commutativity Theorem}
We now prove that non-diagonal transport matrices $M$ or coordinate-scrambling operations destroy representation layers due to non-commutativity with the ReLU activation function $\sigma(x) = \max(0, x)$.

\begin{theorem}[ReLU Commutativity Theorem]
  Let $x \in \mathbb{R}^d$ represent the activations of a layer. Consider a linear transport map $M \in \mathbb{R}^{d \times d}$ applied to the weight rows, which is equivalent to transforming the pre-activations $y = W x$ into $M y$. 
  \begin{enumerate}
    \item If $M = D$ is a positive diagonal matrix, i.e., $D = \mathrm{diag}(d_1, \dots, d_d)$ with $d_i > 0$ for all $i$, then the transformation strictly commutes with the ReLU activation function $\sigma$:
    \begin{equation}
      \sigma(D y) = D \sigma(y) \quad \forall y \in \mathbb{R}^d
    \end{equation}
    \item If $M$ is a non-diagonal matrix, or if $M$ represents a coordinate-scrambling permutation (such as the sorting step in WCPR), then $M$ does not commute with $\sigma$ in general:
    \begin{equation}
      \sigma(M y) \neq M \sigma(y)
    \end{equation}
  \end{enumerate}
\end{theorem}

\begin{proof}
  \textbf{Part 1 (Positive Diagonal Scaling):}
  Let $D = \mathrm{diag}(d_1, \dots, d_d)$ with $d_i > 0$. For any $y \in \mathbb{R}^d$, the $i$-th coordinate of $\sigma(D y)$ is:
  \begin{equation}
    \sigma(D y)_i = \max(0, d_i y_i)
  \end{equation}
  Since $d_i > 0$, we can factor $d_i$ out of the maximum:
  \begin{equation}
    \max(0, d_i y_i) = d_i \max(0, y_i) = d_i \sigma(y_i)
  \end{equation}
  Thus, $\sigma(D y)_i = (D \sigma(y))_i$ for all $i$, which proves $\sigma(D y) = D \sigma(y)$.
  
  \textbf{Part 2 (Non-Diagonal and Permutations):}
  Suppose $M$ is non-diagonal, so there exists $M_{ij} \neq 0$ for $i \neq j$. Consider $y = [1, -1]^T$ and $M = \begin{bmatrix} 1 & 1 \\ 0 & 1 \end{bmatrix}$. Then:
  \begin{equation}
    M y = [0, -1]^T \implies \sigma(M y) = [0, 0]^T
  \end{equation}
  However, running $\sigma(y) = [1, 0]^T$, we get:
  \begin{equation}
    M \sigma(y) = [1, 0]^T \implies \sigma(M y) \neq M \sigma(y)
  \end{equation}
  Similarly, for any coordinate-scrambling permutation matrix $P$, $P$ does not commute with $\sigma$ unless the permutation maps active coordinates onto active coordinates and inactive onto inactive, which is violated for general activations $y$. Thus, non-diagonal transformations and coordinate-scrambling destroy the piecewise linear boundaries established by ReLU, causing catastrophic representation collapse.
\end{proof}

\subsection{Diagonal Bures-Wasserstein Parameter Alignment (D-BWPA)}
Based on Theorem 3.1, we restrict the covariance matrices $\Sigma_{\text{merged}}$ and $\Sigma_{e, k}$ to be diagonal, which corresponds to modeling coordinate-wise variances. For diagonal covariances, the Bures-Wasserstein distance simplifies to the standard Wasserstein distance between independent marginal Gaussian distributions.

For a weight matrix $W \in \mathbb{R}^{d_{\text{out}} \times d_{\text{in}}}$, we compute the row-wise mean and standard deviation:
\begin{align}
  \mu_{\text{merged}, i} &= \frac{1}{d_{\text{in}}} \sum_{j=1}^{d_{\text{in}}} W_{\text{WA}, i, j} \\
  \sigma_{\text{merged}, i} &= \sqrt{\frac{1}{d_{\text{in}}} \sum_{j=1}^{d_{\text{in}}} (W_{\text{WA}, i, j} - \mu_{\text{merged}, i})^2}
\end{align}
For each expert $k$, we compute its row-wise mean $\mu_{e, k, i}$ and standard deviation $\sigma_{e, k, i}$ analogously. The target mean and standard deviation are the averages of the experts' statistics (representing the exact 1D Wasserstein barycenter):
\begin{align}
  \mu_{\text{target}, i} &= \frac{1}{K} \sum_{k=1}^K \mu_{e, k, i} \\
  \sigma_{\text{target}, i} &= \frac{1}{K} \sum_{k=1}^K \sigma_{e, k, i}
\end{align}
The Diagonal Bures-Wasserstein transport map is then given coordinate-wise by:
\begin{equation}
\begin{split}
  T&(W_{\text{WA}})_{i, j} = \mu_{\text{target}, i} \\
  &+ \frac{\sigma_{\text{target}, i}}{\sigma_{\text{merged}, i} + \epsilon} (W_{\text{WA}, i, j} - \mu_{\text{merged}, i})
\end{split}
\end{equation}

Our exact procedure is formalized in Algorithm~\ref{alg:dbwpa}.

\begin{algorithm}[h]
\caption{Diagonal Bures-Wasserstein Parameter Alignment (D-BWPA)}
\label{alg:dbwpa}
\begin{footnotesize}
\begin{algorithmic}[1]
\REQUIRE Expert weights $\{W_{e, k}\}_{k=1}^K$ where $W_{e, k} \in \mathbb{R}^{d_{\text{out}} \times d_{\text{in}}}$, merged weights $W_{\text{WA}} \in \mathbb{R}^{d_{\text{out}} \times d_{\text{in}}}$, regularization $\epsilon > 0$.
\ENSURE Aligned weights $W_{\text{aligned}} \in \mathbb{R}^{d_{\text{out}} \times d_{\text{in}}}$.
\FOR{$i = 1$ \TO $d_{\text{out}}$}
  \STATE Compute merged mean: $\mu_{\text{merged}, i} \gets \frac{1}{d_{\text{in}}} \sum_{j=1}^{d_{\text{in}}} W_{\text{WA}, i, j}$
  \STATE Compute merged standard deviation: $\sigma_{\text{merged}, i} \gets \sqrt{\frac{1}{d_{\text{in}}} \sum_{j=1}^{d_{\text{in}}} (W_{\text{WA}, i, j} - \mu_{\text{merged}, i})^2}$
  \FOR{$k = 1$ \TO $K$}
    \STATE Compute expert mean: $\mu_{e, k, i} \gets \frac{1}{d_{\text{in}}} \sum_{j=1}^{d_{\text{in}}} W_{e, k, i, j}$
    \STATE Compute expert standard deviation: $\sigma_{e, k, i} \gets \sqrt{\frac{1}{d_{\text{in}}} \sum_{j=1}^{d_{\text{in}}} (W_{e, k, i, j} - \mu_{e, k, i})^2}$
  \ENDFOR
  \STATE Compute target mean barycenter: $\mu_{\text{target}, i} \gets \frac{1}{K} \sum_{k=1}^K \mu_{e, k, i}$
  \STATE Compute target standard deviation barycenter: $\sigma_{\text{target}, i} \gets \frac{1}{K} \sum_{k=1}^K \sigma_{e, k, i}$
  \FOR{$j = 1$ \TO $d_{\text{in}}$}
    \STATE Aligned element: $W_{\text{aligned}, i, j} \gets \mu_{\text{target}, i} + \frac{\sigma_{\text{target}, i}}{\sigma_{\text{merged}, i} + \epsilon} (W_{\text{WA}, i, j} - \mu_{\text{merged}, i})$
  \ENDFOR
\ENDFOR
\end{algorithmic}
\end{footnotesize}
\end{algorithm}

This formulation has substantial benefits over full-covariance BWPA: (i) \textbf{Strict ReLU-Commutativity:} The scaling factor $\frac{\sigma_{\text{target}, i}}{\sigma_{\text{merged}, i} + \epsilon} > 0$ strictly commutes with ReLU, preventing representation collapse on MLPs. (ii) \textbf{Numerical Stability:} Unlike full-covariance alignment, D-BWPA requires no matrix inversions or eigen-decompositions. It is robust to rank-deficient weight matrices where $d_{\text{in}} < d_{\text{out}}$. (iii) \textbf{Computational Efficiency:} The time complexity is $O(d_{\text{out}} d_{\text{in}})$—which is linear in the parameter count—compared to the cubic $O(d_{\text{out}}^3)$ complexity of full-covariance alignment.

\subsection{Computational Complexity Analysis}
To establish the practical viability of our proposed Diagonal Bures-Wasserstein Parameter Alignment (D-BWPA), we formalize its computational complexity relative to sorting-based 1D optimal transport (WCPR) and full-covariance Bures-Wasserstein Parameter Alignment (Full-BWPA).

\begin{proposition}[Computational Complexity]
Let $W \in \mathbb{R}^{d_{\text{out}} \times d_{\text{in}}}$ be a weight matrix, and let $K$ be the number of expert models. 
\begin{enumerate}
  \item Sorting-based WCPR has a time complexity of $\mathcal{O}(K d_{\text{out}} d_{\text{in}} \log d_{\text{in}})$.
  \item Full-covariance BWPA, using fixed-point iteration for the barycenter of $d_{\text{out}} \times d_{\text{out}}$ matrices, has a time complexity of $\mathcal{O}(T d_{\text{out}}^3 + K d_{\text{out}} d_{\text{in}})$, where $T$ is the number of fixed-point iterations.
  \item Our proposed D-BWPA has a time complexity of $\mathcal{O}(K d_{\text{out}} d_{\text{in}})$, which is strictly linear in the number of network parameters.
\end{enumerate}
\end{proposition}

\begin{proof}
1. WCPR operates channel-by-channel (row-by-row). For each of the $d_{\text{out}}$ rows, WCPR performs sorting on the $d_{\text{in}}$ flat elements of the merged weights and each of the $K$ expert weights. The complexity of sorting $d_{\text{in}}$ elements is $\mathcal{O}(d_{\text{in}} \log d_{\text{in}})$. Done across $K$ experts and $d_{\text{out}}$ rows, the total complexity is $\mathcal{O}(K d_{\text{out}} d_{\text{in}} \log d_{\text{in}})$.
2. Full-covariance BWPA computes the $d_{\text{out}} \times d_{\text{out}}$ covariance matrix for each of the $K$ experts and the merged weights, requiring $\mathcal{O}(K d_{\text{out}}^2 d_{\text{in}} + d_{\text{out}}^2 d_{\text{in}})$ operations. It then solves the Bures-Wasserstein fixed-point iteration on $d_{\text{out}} \times d_{\text{out}}$ matrices, where each iteration requires matrix square roots and matrix multiplications of cost $\mathcal{O}(K d_{\text{out}}^3)$. Assuming $T$ iterations and since $d_{\text{in}} \ge d_{\text{out}}$ in typical deep learning layers, the total complexity is $\mathcal{O}(T d_{\text{out}}^3 + K d_{\text{out}} d_{\text{in}})$.
3. D-BWPA computes the coordinate-wise mean and variance of $d_{\text{in}}$ elements for each of the $d_{\text{out}}$ rows. This requires $\mathcal{O}(d_{\text{in}})$ operations per row. Done for $K$ experts and the merged weight matrix, the total statistics computation takes $\mathcal{O}(K d_{\text{out}} d_{\text{in}})$ operations. The subsequent element-wise transport map is computed in $\mathcal{O}(d_{\text{out}} d_{\text{in}})$ time. Thus, the total complexity of D-BWPA is $\mathcal{O}(K d_{\text{out}} d_{\text{in}})$, which is linear in the parameter size and independent of logarithmic sorting or cubic matrix inversion overheads.
\end{proof}

\subsection{Quantization Stability and Range Guarantees}
Next, we analyze the sensitivity of model merging algorithms to post-training quantization. Post-training 8-bit quantization scales and clamps a continuous parameter $w$ to a discrete grid:
\begin{equation}
  q(w) = \mathrm{clamp}\left( \left\lfloor \frac{w}{S} \right\rceil, -128, 127 \right) \cdot S
\end{equation}
where the scale factor $S$ is defined by the peak dynamic range:
\begin{equation}
  S = \frac{\max |w|}{127}
\end{equation}
Uniform quantization introduces a noise term $e(w) = q(w) - w$, bounded by $|e(w)| \le S/2$.

\begin{proposition}[Quantization Noise Stability]
Let $T(W)$ be the transported weight matrix.
\begin{enumerate}
  \item For Diagonal Bures-Wasserstein Parameter Alignment (D-BWPA), the transported weights are bounded, and their maximum absolute value is strictly bounded by the expert statistics, preventing dynamic range expansion:
  \begin{equation}
  \begin{split}
    |T&(W_{\mathrm{WA}})_{i,j}| \le |\mu_{\mathrm{target}, i}| \\
    &+ \frac{\sigma_{\mathrm{target}, i}}{\sigma_{\mathrm{merged}, i} + \epsilon} |W_{\mathrm{WA}, i,j} - \mu_{\mathrm{merged}, i}|
  \end{split}
  \end{equation}
  \item For sorting-based WCPR, if the probability density of the merged weight rows has thin tails (e.g. Gaussian or sub-Gaussian), the inverse cumulative distribution function $F_{\mathrm{target}}^{-1}(p)$ in WCPR maps outliers (where $p \to 0$ or $p \to 1$) to highly magnified values. This arbitrarily inflates the dynamic range $\max |T(W_{\mathrm{WCPR}})|$, thereby increasing the step size $S$ and amplifying the quantization noise $|e(w)|$ across the entire parameter tensor.
\end{enumerate}
\end{proposition}

\begin{proof}
1. In D-BWPA, the scaling factor is a smooth, regularized ratio of standard deviations: $c_i = \sigma_{\mathrm{target}, i} / (\sigma_{\mathrm{merged}, i} + \epsilon)$. Since $\sigma_{\mathrm{target}}$ is the average of the experts' standard deviations, $c_i \approx 1$ under standard fine-tuning from a shared progenitor. The range of the transported weights remains extremely close to the expert weights' original range, maintaining a small step size $S_{\mathrm{D-BWPA}}$ and low quantization noise.
2. In WCPR, the quantile-matching transport is $T_{\mathrm{WCPR}}(w) = F_{\mathrm{target}}^{-1}(F_{\mathrm{merged}}(w))$. For a normal distribution $\mathcal{N}(\mu, \sigma^2)$, $F(w)$ approaches 1 as $w \to \infty$. In the tail of the merged weights, a small numerical perturbation or outlier weight $w_{\mathrm{out}}$ can yield $F_{\mathrm{merged}}(w_{\mathrm{out}}) = 1 - \delta$. Under the inverse CDF mapping, $F_{\mathrm{target}}^{-1}(1 - \delta) = \mu_{\mathrm{target}} + \sigma_{\mathrm{target}} \Phi^{-1}(1 - \delta)$, where $\Phi^{-1}$ is the probit function. As $\delta \to 0$, $\Phi^{-1}(1 - \delta) \to \infty$ at an asymptotic rate of $\mathcal{O}(\sqrt{-\log \delta})$. This quantile stretching in the tails creates highly isolated, extreme-value outliers in the merged tensor. This increases the peak value $\max |T(W_{\mathrm{WCPR}})|$ dramatically. Since the scale factor $S = \max |T(W)| / 127$ depends on the single maximum absolute value, this outlier-induced peak inflates $S$ for the entire layer, causing severe quantization rounding errors and representation collapse for the remaining $99.9\%$ of the weights.
\end{proof}

\subsection{Progenitor-Centric Distance and Stability Bounds}
A critical requirement for model merging algorithms is parameter-space stability: the merged model must remain within the low-loss basin surrounding the progenitor model. If the alignment transport map distorts weights excessively, the model may leave this basin, causing optimization collapse. We now prove that our proposed D-BWPA strictly preserves progenitor-centric parameter distance.

\begin{proposition}[Distance to Progenitor]
Let $w_i, w_{\mathrm{prog}, i}, w_{\mathrm{aligned}, i} \in \mathbb{R}^{d_{\mathrm{in}}}$ represent the $i$-th row of $W_{\mathrm{WA}}$, $W_{\mathrm{prog}}$, and the aligned weights $W_{\mathrm{aligned}}$ under D-BWPA for any regularization parameter $\epsilon \ge 0$, respectively. Let $\mu_{\mathrm{target}, i}$, $\mu_{\mathrm{prog}, i}$ and $\sigma_{\mathrm{target}, i}$, $\sigma_{\mathrm{prog}, i}$ be the row-wise means and standard deviations. Let $\rho_i \in [-1, 1]$ denote the Pearson correlation coefficient between the zero-mean components of the merged weights and the progenitor weights. Then:
\begin{equation}
\begin{split}
  &\|w_{\mathrm{aligned}, i} - w_{\mathrm{prog}, i}\|^2 \\
  &= d_{\mathrm{in}} (\mu_{\mathrm{target}, i} - \mu_{\mathrm{prog}, i})^2 \\
  &\quad + d_{\mathrm{in}} \left( \sigma_{\mathrm{target}, i} \frac{\sigma_{\mathrm{merged}, i}}{\sigma_{\mathrm{merged}, i} + \epsilon} - \sigma_{\mathrm{prog}, i} \right)^2 \\
  &\quad + 2 d_{\mathrm{in}} \sigma_{\mathrm{target}, i} \sigma_{\mathrm{prog}, i} \frac{\sigma_{\mathrm{merged}, i}}{\sigma_{\mathrm{merged}, i} + \epsilon} (1 - \rho_i)
\end{split}
\end{equation}
\end{proposition}

\begin{proof}
Let $\mathbf{1} \in \mathbb{R}^{d_{\mathrm{in}}}$ be the vector of all ones. Let $\tilde{w}_i = w_i - \mu_{\mathrm{merged}, i} \mathbf{1}$ and $\tilde{w}_{\mathrm{prog}, i} = w_{\mathrm{prog}, i} - \mu_{\mathrm{prog}, i} \mathbf{1}$ be the zero-mean components of the weights. By definition of standard deviation, we have $\|\tilde{w}_i\|^2 = d_{\mathrm{in}} \sigma_{\mathrm{merged}, i}^2$ and $\|\tilde{w}_{\mathrm{prog}, i}\|^2 = d_{\mathrm{in}} \sigma_{\mathrm{prog}, i}^2$.
Using D-BWPA for any $\epsilon \ge 0$, we have:
\begin{equation}
  w_{\mathrm{aligned}, i} = \mu_{\mathrm{target}, i} \mathbf{1} + \frac{\sigma_{\mathrm{target}, i}}{\sigma_{\mathrm{merged}, i} + \epsilon} \tilde{w}_i
\end{equation}
The difference from the progenitor row is:
\begin{equation}
\begin{split}
  w_{\mathrm{aligned}, i} &- w_{\mathrm{prog}, i} = (\mu_{\mathrm{target}, i} - \mu_{\mathrm{prog}, i}) \mathbf{1} \\
  &+ \left( \frac{\sigma_{\mathrm{target}, i}}{\sigma_{\mathrm{merged}, i} + \epsilon} \tilde{w}_i - \tilde{w}_{\mathrm{prog}, i} \right)
\end{split}
\end{equation}
Note that $\mathbf{1}$ is orthogonal to both zero-mean vectors $\tilde{w}_i$ and $\tilde{w}_{\mathrm{prog}, i}$. Therefore, the squared $L_2$ norm decomposes exactly into:
\begin{equation}
\begin{split}
  &\|w_{\mathrm{aligned}, i} - w_{\mathrm{prog}, i}\|^2 \\
  &= \|(\mu_{\mathrm{target}, i} - \mu_{\mathrm{prog}, i}) \mathbf{1}\|^2 \\
  &+ \left\| \frac{\sigma_{\mathrm{target}, i}}{\sigma_{\mathrm{merged}, i} + \epsilon} \tilde{w}_i - \tilde{w}_{\mathrm{prog}, i} \right\|^2 \\
  &= d_{\mathrm{in}} (\mu_{\mathrm{target}, i} - \mu_{\mathrm{prog}, i})^2 \\
  &+ \left( \frac{\sigma_{\mathrm{target}, i}}{\sigma_{\mathrm{merged}, i} + \epsilon} \right)^2 \|\tilde{w}_i\|^2 + \|\tilde{w}_{\mathrm{prog}, i}\|^2 \\
  &- 2 \frac{\sigma_{\mathrm{target}, i}}{\sigma_{\mathrm{merged}, i} + \epsilon} \langle \tilde{w}_i, \tilde{w}_{\mathrm{prog}, i} \rangle
\end{split}
\end{equation}
Substituting $\|\tilde{w}_i\|^2 = d_{\mathrm{in}} \sigma_{\mathrm{merged}, i}^2$ and $\|\tilde{w}_{\mathrm{prog}, i}\|^2 = d_{\mathrm{in}} \sigma_{\mathrm{prog}, i}^2$:
\begin{equation}
\begin{split}
  &\|w_{\mathrm{aligned}, i} - w_{\mathrm{prog}, i}\|^2 \\
  &= d_{\mathrm{in}} (\mu_{\mathrm{target}, i} - \mu_{\mathrm{prog}, i})^2 \\
  &+ d_{\mathrm{in}} \sigma_{\mathrm{target}, i}^2 \frac{\sigma_{\mathrm{merged}, i}^2}{(\sigma_{\mathrm{merged}, i} + \epsilon)^2} + d_{\mathrm{in}} \sigma_{\mathrm{prog}, i}^2 \\
  &- 2 \frac{\sigma_{\mathrm{target}, i}}{\sigma_{\mathrm{merged}, i} + \epsilon} \langle \tilde{w}_i, \tilde{w}_{\mathrm{prog}, i} \rangle
\end{split}
\end{equation}
Recall the correlation coefficient:
\begin{equation}
  \rho_i = \frac{\langle \tilde{w}_i, \tilde{w}_{\mathrm{prog}, i} \rangle}{d_{\mathrm{in}} \sigma_{\mathrm{merged}, i} \sigma_{\mathrm{prog}, i}}
\end{equation}
This gives $\langle \tilde{w}_i, \tilde{w}_{\mathrm{prog}, i} \rangle = d_{\mathrm{in}} \sigma_{\mathrm{merged}, i} \sigma_{\mathrm{prog}, i} \rho_i$. Substituting this in:
\begin{equation}
\begin{split}
  &\|w_{\mathrm{aligned}, i} - w_{\mathrm{prog}, i}\|^2 \\
  &= d_{\mathrm{in}} (\mu_{\mathrm{target}, i} - \mu_{\mathrm{prog}, i})^2 \\
  &\quad + d_{\mathrm{in}} \sigma_{\mathrm{target}, i}^2 \frac{\sigma_{\mathrm{merged}, i}^2}{(\sigma_{\mathrm{merged}, i} + \epsilon)^2} + d_{\mathrm{in}} \sigma_{\mathrm{prog}, i}^2 \\
  &\quad - 2 d_{\mathrm{in}} \sigma_{\mathrm{target}, i} \sigma_{\mathrm{prog}, i} \rho_i \frac{\sigma_{\mathrm{merged}, i}}{\sigma_{\mathrm{merged}, i} + \epsilon} \\
  &= d_{\mathrm{in}} (\mu_{\mathrm{target}, i} - \mu_{\mathrm{prog}, i})^2 \\
  &\quad + d_{\mathrm{in}} \left( \sigma_{\mathrm{target}, i} \frac{\sigma_{\mathrm{merged}, i}}{\sigma_{\mathrm{merged}, i} + \epsilon} - \sigma_{\mathrm{prog}, i} \right)^2 \\
  &\quad + 2 d_{\mathrm{in}} \sigma_{\mathrm{target}, i} \sigma_{\mathrm{prog}, i} \frac{\sigma_{\mathrm{merged}, i}}{\sigma_{\mathrm{merged}, i} + \epsilon} (1 - \rho_i)
\end{split}
\end{equation}
which completes the proof.
\end{proof}

This elegant relation reveals that parameter-space distance from the aligned model to the progenitor is strictly bounded. Under standard fine-tuning from a shared progenitor, the expert weights remain highly correlated with the progenitor ($\rho_i \approx 1$) and have similar means and variances ($\mu_{\mathrm{target}, i} \approx \mu_{\mathrm{prog}, i}$ and $\sigma_{\mathrm{target}, i} \approx \sigma_{\mathrm{prog}, i}$), ensuring that D-BWPA resides stably within the progenitor's local basin of convergence.

\begin{table*}[t]
\caption{ResNet-18 Merging Accuracy on MNIST, FashionMNIST, and CIFAR-10.}
\label{tab:resnet}
\vskip 0.05in
\begin{center}
\begin{footnotesize}
\begin{sc}
\setlength{\tabcolsep}{1.5pt}
\begin{tabular}{lcccc}
\toprule
Method & MNIST & Fashion & CIFAR-10 & Average \\
\midrule
\textbf{Oracle (Experts)} & 98.98\% & 92.27\% & 80.84\% & \textbf{90.70\%} \\
\midrule
Weight Averaging (WA) & 38.58\% & 17.02\% & 22.01\% & 25.87\% \\
WA + 8-bit Quantization & 38.52\% & 17.08\% & 22.05\% & 25.88\% \\
WA + DE-BN (Corrected) & 93.03\% & 70.46\% & 41.42\% & \textbf{68.30\%} \\
WA + DE-BN + 8-bit Quant (Corrected) & 91.34\% & 75.25\% & 40.72\% & \textbf{69.10\%} \\
\midrule
Tuned Task Arithmetic (TA) & 40.76\% & 19.49\% & 27.95\% & 29.40\% \\
Tuned TA + 8-bit Quant & 40.43\% & 19.81\% & 27.62\% & 29.29\% \\
Tuned TA + DE-BN & 88.87\% & 71.33\% & 43.63\% & 67.94\% \\
Tuned TA + DE-BN + 8-bit Quant & 91.09\% & 75.26\% & 43.11\% & 69.82\% \\
\midrule
Isotropic Parameter Resonance (U-IPR) & 9.80\% & 10.00\% & 10.00\% & 9.93\% \\
Spectral Parameter Resonance (S-IPR) & 72.73\% & 52.68\% & 31.11\% & 52.17\% \\
S-IPR + 8-bit Quantization & 73.29\% & 52.63\% & 31.32\% & 52.41\% \\
Wasserstein Calibration (WCPR) & 48.88\% & 20.52\% & 25.08\% & 31.49\% \\
WCPR + 8-bit Quantization & 48.32\% & 20.66\% & 25.19\% & 31.39\% \\
\midrule
\textbf{D-BWPA (Ours)} & 48.86\% & 21.19\% & 25.33\% & \textbf{31.79\%} \\
\textbf{D-BWPA (Ours) + 8-bit Quant} & 49.60\% & 21.15\% & 25.44\% & \textbf{32.06\%} \\
\textbf{D-BWPA (Ours) + DE-BN} & 92.83\% & 71.94\% & 38.28\% & \textbf{67.68\%} \\
\textbf{D-BWPA (Ours) + DE-BN + 8-bit Quant} & 92.06\% & 68.26\% & 43.12\% & \textbf{67.81\%} \\
\midrule
Full-BWPA (Ablation) & 31.68\% & 13.94\% & 12.13\% & 19.25\% \\
Full-BWPA (Ablation) + 8-bit Quant & 32.17\% & 14.35\% & 12.26\% & 19.59\% \\
Full-BWPA + DE-BN & 69.28\% & 46.91\% & 27.04\% & 47.74\% \\
Full-BWPA + DE-BN + 8-bit Quant & 67.41\% & 40.37\% & 27.29\% & 45.02\% \\
\bottomrule
\end{tabular}
\end{sc}
\end{footnotesize}
\end{center}
\vskip -0.15in
\end{table*}

\section{Algorithmic Resolution of BatchNorm Calibration}
\label{sec:bn_cal}

When merging convolutional networks like ResNet-18, the alignment of Batch Normalization statistics is critical to prevent activation mismatch. In data-efficient BN calibration (DE-BN), we update the running mean and variance using a small calibration set.

\subsection{The Momentum Mismatch Bug}
In PyTorch, the running statistics $\mu_{\text{running}}$ and $\sigma^2_{\text{running}}$ of a BatchNorm layer are updated during training using exponential moving average (EMA) with momentum $\beta \in [0, 1]$ (default $\beta = 0.1$):
\begin{align}
  \mu_{\text{running}}^{(t+1)} &= (1 - \beta) \mu_{\text{running}}^{(t)} + \beta \mu_{\text{batch}} \\
  \sigma_{\text{running}}^{2(t+1)} &= (1 - \beta) \sigma_{\text{running}}^{2(t)} + \beta \sigma_{\text{batch}}^2
\end{align}
In standard DE-BN baselines, the running statistics are reset to $0$ and $1$, and a single batch of calibration data is forwarded.
However, with a single batch forward and $\beta = 0.1$, the running statistics become:
\begin{align}
  \mu_{\text{running}}^{(1)} &= 0.9 \times 0 + 0.1 \mu_{\text{batch}} = 0.1 \mu_{\text{batch}} \\
  \sigma_{\text{running}}^{2(1)} &= 0.9 \times 1 + 0.1 \sigma_{\text{batch}}^2 = 0.9 + 0.1 \sigma_{\text{batch}}^2
\end{align}
Thus, the running statistics are still 90\% stuck at their default uncalibrated initialization values! This causes severe representation collapse when the model is switched back to evaluation mode, explaining why prior WA+DE-BN baselines achieved only 10.03% (random guessing) on ResNet-18.

\subsection{The Momentum-Corrected DE-BN Algorithm}
To resolve this, we force $\beta = 1.0$ during the calibration forward pass. This forces the running statistics to exactly capture the calibration batch statistics:
\begin{align}
  \mu_{\text{running}}^{(1)} &= \mu_{\text{batch}} \\
  \sigma_{\text{running}}^{2(1)} &= \sigma_{\text{batch}}^2
\end{align}

Our complete, momentum-corrected DE-BN procedure is formalized in Algorithm~\ref{alg:debn}.

\begin{algorithm}[h]
\caption{Momentum-Corrected Data-Efficient BatchNorm Calibration}
\label{alg:debn}
\begin{footnotesize}
\begin{algorithmic}[1]
\REQUIRE Merged model $f_{\theta}$ with Batch Normalization layers $\{\mathrm{BN}_l\}$, calibration dataset $\mathcal{D}_{\text{cal}}$ (typically $8$ to $32$ samples), default momentum values $\{\beta_l\}$.
\ENSURE Calibrated model $f_{\theta}$ with adapted running statistics.
\FOR{each Batch Normalization layer $\mathrm{BN}_l$}
  \STATE Reset running mean: $\mu_{\text{running}, l} \gets 0$
  \STATE Reset running variance: $\sigma^2_{\text{running}, l} \gets 1$
  \STATE Save default momentum and set calibration momentum: $\beta_l^{\text{temp}} \gets \beta_l, \quad \beta_l \gets 1.0$
\ENDFOR
\STATE Set model to training mode: $f_{\dots}.\mathrm{train}()$ (enables statistics tracking)
\STATE Forward pass calibration batch: $X \gets \mathcal{D}_{\text{cal}} \implies f_{\theta}(X)$
\FOR{each Batch Normalization layer $\mathrm{BN}_l$}
  \STATE Restore default momentum: $\beta_l \gets \beta_l^{\text{temp}}$
\ENDFOR
\STATE Set model to evaluation mode: $f_{\theta}.\mathrm{eval}()$
\end{algorithmic}
\end{footnotesize}
\end{algorithm}

Once the single-batch forward pass is completed, we restore the original momentum values of the BN layers. This elegant algorithmic change resolves the calibration issue, as we demonstrate in our experimental validation.

\section{Experiments and Results}
\label{sec:experiments}
We evaluated our proposed D-BWPA and the momentum-corrected DE-BN algorithm across ResNet-18 and MLP architectures. Task-specific experts were trained on MNIST, FashionMNIST, and CIFAR-10, fine-tuned from a shared progenitor model for 5 epochs. We evaluated merging methods under standard precision and post-training 8-bit uniform quantization.

\begin{figure*}[ht]
  \centering
  \includegraphics[width=0.38\linewidth]{merging_comparison}
  \caption{Comparison of average model merging accuracy across different settings. Panel (a) illustrates performance on MLP architectures under standard FP32 and 8-bit post-training quantization, showcasing the representation collapse of Full-BWPA and the resilience of D-BWPA. Panel (b) illustrates ResNet-18 merging accuracy under 8-bit quantization with and without our momentum-corrected data-efficient BatchNorm calibration (DE-BN), showing a massive accuracy boost for all valid methods.}
  \label{fig:merging_comparison}
  \vspace{-8pt}
\end{figure*}

\subsection{Quantitative Results on ResNet-18}
Our results on ResNet-18 are summarized in Table~\ref{tab:resnet}. 

Prior work reported a catastrophic 10.03\% average accuracy for WA + DE-BN due to the momentum bug. Our simple, theoretically sound 1-line momentum correction resolves this, immediately restoring WA + DE-BN to \textbf{68.30\%} and our D-BWPA + DE-BN to \textbf{67.68\%} in our evaluations. 

Furthermore, our D-BWPA alone (without DE-BN) achieves \textbf{31.79\%} (and \textbf{32.06\%} under 8-bit quantization), outperforming the competitive Wasserstein baseline WCPR (31.49\%). This indicates that our coordinate-preserving, linear-covariant transport provides a more robust alignment of model weights.

\textbf{Full-Covariance Ablation (Full-BWPA):} As shown in Table~\ref{tab:resnet}, removing the diagonal covariance restriction (Full-BWPA) causes severe performance collapse. Without DE-BN, Full-BWPA drops to \textbf{19.25\%} (versus \textbf{31.79\%} for D-BWPA). Even with DE-BN, Full-BWPA + DE-BN recovers to only \textbf{47.74\%}, compared to \textbf{67.68\%} for D-BWPA + DE-BN. This confirms that non-diagonal coordinate scrambling severely corrupts layer-wise representations, which downstream BatchNorm calibration cannot fully repair.

\subsection{Quantitative Results on MLP}
Our results on MLP architectures are shown in Table~\ref{tab:mlp}.

\begin{table}[h]
\caption{MLP Merging Accuracy (No BatchNorm).}
\label{tab:mlp}
\vskip 0.05in
\begin{center}
\begin{footnotesize}
\begin{sc}
\setlength{\tabcolsep}{1.0pt}
\begin{tabular}{lcccc}
\toprule
Method & MNIST & Fashion & CIFAR-10 & Avg \\
\midrule
\textbf{Oracle (Experts)} & 97.27\% & 87.24\% & 54.29\% & \textbf{79.60\%} \\
\midrule
WA (Baseline) & 47.74\% & 42.32\% & 24.23\% & 38.10\% \\
WA + 8-bit Q & 47.73\% & 42.35\% & 24.21\% & 38.10\% \\
Tuned TA & 50.22\% & 44.67\% & 25.41\% & 40.10\% \\
Tuned TA + 8-bit Q & 50.27\% & 44.68\% & 25.42\% & 40.12\% \\
U-IPR & 48.61\% & 43.21\% & 24.65\% & 38.82\% \\
S-IPR & 41.33\% & 39.77\% & 24.39\% & 35.16\% \\
WCPR & 52.47\% & 48.79\% & 24.63\% & \textbf{41.96\%} \\
WCPR + 8-bit Q & 52.48\% & 48.84\% & 24.63\% & \textbf{41.98\%} \\
\midrule
\textbf{D-BWPA (Ours)} & 50.12\% & 46.70\% & 24.67\% & \textbf{40.50\%} \\
\textbf{D-BWPA (Ours) + Q} & 50.17\% & 46.77\% & 24.64\% & \textbf{40.53\%} \\
\midrule
Full-BWPA (Abl) & 25.58\% & 12.68\% & 12.21\% & 16.82\% \\
Full-BWPA (Abl) + Q & 25.35\% & 12.67\% & 12.21\% & 16.74\% \\
\bottomrule
\end{tabular}
\end{sc}
\end{footnotesize}
\end{center}
\vskip -0.18in
\vspace{-5pt}
\end{table}

While Spectral Parameter Resonance (S-IPR) collapses to 35.16\% (significantly worse than base Weight Averaging) due to representation scrambling, our proposed D-BWPA achieves an excellent \textbf{40.50\%} average accuracy, outperforming both WA and Tuned TA. Under 8-bit post-training quantization, D-BWPA demonstrates absolute robustness, actually improving to \textbf{40.53\%}. This strongly supports Theorem 3.1, proving that our diagonal formulation successfully preserves the coordinate-wise representations necessary for ReLU-based networks.

\textbf{Failure of Full-Covariance BWPA:} This is most clearly illustrated by the Full-BWPA ablation. Under full-covariance transport, the MLP experiences catastrophic collapse (16.82\% average accuracy). Because Full-BWPA uses a non-diagonal transport matrix, it mixes coordinates across activation channels, violating ReLU-commutativity (Theorem 3.1) and destroying linear activation boundaries. Conversely, WCPR achieves 41.96\% average accuracy because, although it sorts weight values, it preserves their coordinate indices. This highlights the necessity of ReLU-commutative diagonal transformations (such as positive diagonal scaling in D-BWPA) for stable, data-free parameter alignment in feedforward architectures.

\subsection{Sensitivity Analysis, Limitations, and Impact}
\label{sec:discussion}

\textbf{Regularization and Numerical Stability.}
Our proposed D-BWPA algorithm depends on a small numerical regularization parameter $\epsilon > 0$ to avoid division by zero when standard deviations approach zero. We observed that the performance of D-BWPA is highly stable across a wide range of values $\epsilon \in [10^{-7}, 10^{-4}]$. When $\epsilon$ becomes too large ($\ge 10^{-2}$), the scaling factor $\frac{\sigma_{\text{target}, i}}{\sigma_{\text{merged}, i} + \epsilon}$ shrinks towards zero, leading to under-scaled parameter magnitudes. Conversely, when $\epsilon$ is too small ($\le 10^{-9}$), the numerical division becomes unstable on coordinate channels with extremely sparse activations. Thus, $\epsilon = 10^{-6}$ serves as a robust, scale-invariant default across all models and datasets.

\textbf{Limitations of Diagonal Alignment.}
While D-BWPA is provably ReLU-commutative and computationally efficient, its primary limitation is that it models only diagonal covariance structures. By neglecting the off-diagonal covariance entries (which capture inter-channel correlation), D-BWPA assumes independent marginal distributions for weight rows. This independence assumption restricts the complexity of representation alignment. However, as demonstrated by our Full-BWPA ablation, any attempt to incorporate off-diagonal elements through non-diagonal transport matrices violates ReLU-commutativity and scrambles representation structures. Therefore, diagonal alignment represents a fundamental, theoretically bounded trade-off between alignment complexity and structural preservation.

\textbf{Ethical Considerations and Energy Footprints.}
Serving specialized model experts imposes substantial computational and memory footprints. Traditional multi-task consolidation requires joint training on multiple datasets, consuming immense GPU energy and producing a significant carbon footprint. In contrast, data-free model merging integrates expert networks within seconds on a single CPU. By resolving the representation collapse of weight averaging without retraining, our proposed D-BWPA and momentum-corrected DE-BN methods promote highly sustainable, energy-efficient deep learning workflows.

\section{Conclusion and Future Work}
In this work, we provided a rigorous theoretical deconstruction of optimal transport and calibration in parameter-space model merging. We proved the \emph{ReLU Commutativity Theorem}, showing that positive diagonal scaling strictly commutes with ReLU, and formulated \emph{Diagonal Bures-Wasserstein Parameter Alignment (D-BWPA)} to avoid the collapse of non-diagonal alignment. Finally, we resolved a critical momentum-update bug in standard BatchNorm calibration (DE-BN), immediately restoring ResNet-18 merging accuracy. Future work will extend ReLU-commutative optimal transport to Transformer architectures.

\clearpage
\bibliography{submission}
\bibliographystyle{icml2026}

\clearpage
\appendix
\onecolumn
\section{Extended Theoretical Derivations}
\label{sec:appendix_theory}

We provide the complete derivations and proofs for the theoretical statements made in Section 3 of the main paper.

\subsection{Equivalence of Bures-Wasserstein and Marginal Wasserstein-2 under Diagonal Covariances}
We derive the exact connection between the multivariate Bures-Wasserstein metric and univariate Wasserstein-2 distances under the assumption of diagonal covariance matrices.

\begin{proposition}[Diagonal Bures-Wasserstein Equivalence]
Let $P_1 = \mathcal{N}(\mu_1, \Sigma_1)$ and $P_2 = \mathcal{N}(\mu_2, \Sigma_2)$ be multivariate Gaussian distributions in $\mathbb{R}^d$, where the covariance matrices $\Sigma_1$ and $\Sigma_2$ are diagonal:
\begin{equation}
  \Sigma_1 = \mathrm{diag}(\sigma_{1, 1}^2, \dots, \sigma_{1, d}^2), \quad \Sigma_2 = \mathrm{diag}(\sigma_{2, 1}^2, \dots, \sigma_{2, d}^2)
\end{equation}
Then, the Bures-Wasserstein distance $d_{\mathrm{BW}}(P_1, P_2)$ simplifies exactly to the root-sum-of-squares of the independent 1D marginal Wasserstein-2 distances between each coordinate:
\begin{equation}
  d_{\mathrm{BW}}^2(P_1, P_2) = \sum_{j=1}^d d_{\mathrm{W}_2}^2\left( \mathcal{N}(\mu_{1, j}, \sigma_{1, j}^2), \, \mathcal{N}(\mu_{2, j}, \sigma_{2, j}^2) \right)
\end{equation}
\end{proposition}

\begin{proof}
The general closed-form Bures-Wasserstein distance between two multivariate Gaussians $P_1$ and $P_2$ is:
\begin{equation}
  d_{\mathrm{BW}}^2(P_1, P_2) = \|\mu_1 - \mu_2\|^2 + \mathrm{Tr}\left( \Sigma_1 + \Sigma_2 - 2 \left( \Sigma_1^{1/2} \Sigma_2 \Sigma_1^{1/2} \right)^{1/2} \right)
\end{equation}
Since $\Sigma_1$ and $\Sigma_2$ are diagonal matrices, they are symmetric positive-definite and commute. Therefore, we can simplify the product term:
\begin{equation}
  \Sigma_1^{1/2} \Sigma_2 \Sigma_1^{1/2} = \Sigma_1 \Sigma_2
\end{equation}
Taking the matrix square root of this diagonal product yields:
\begin{equation}
  \left( \Sigma_1 \Sigma_2 \right)^{1/2} = \mathrm{diag}(\sigma_{1, 1}\sigma_{2, 1}, \dots, \sigma_{1, d}\sigma_{2, d})
\end{equation}
Substituting this back into the trace term:
\begin{equation}
\begin{split}
  \mathrm{Tr}\left( \Sigma_1 + \Sigma_2 - 2 \left( \Sigma_1 \Sigma_2 \right)^{1/2} \right) &= \sum_{j=1}^d \left( \sigma_{1, j}^2 + \sigma_{2, j}^2 - 2 \sigma_{1, j} \sigma_{2, j} \right) \\
  &= \sum_{j=1}^d (\sigma_{1, j} - \sigma_{2, j})^2
\end{split}
\end{equation}
Combining the mean difference and variance difference terms:
\begin{equation}
  d_{\mathrm{BW}}^2(P_1, P_2) = \sum_{j=1}^d \left( (\mu_{1, j} - \mu_{2, j})^2 + (\sigma_{1, j} - \sigma_{2, j})^2 \right)
\end{equation}
Recall that the 1D Wasserstein-2 distance between two univariate Gaussians $\mathcal{N}(\mu_x, \sigma_x^2)$ and $\mathcal{N}(\mu_y, \sigma_y^2)$ is:
\begin{equation}
  d_{\mathrm{W}_2}^2\left( \mathcal{N}(\mu_x, \sigma_x^2), \, \mathcal{N}(\mu_y, \sigma_y^2) \right) = (\mu_x - \mu_y)^2 + (\sigma_x - \sigma_y)^2
\end{equation}
Substituting this 1D relation coordinate-wise completed the proof:
\begin{equation}
  d_{\mathrm{BW}}^2(P_1, P_2) = \sum_{j=1}^d d_{\mathrm{W}_2}^2\left( \mathcal{N}(\mu_{1, j}, \sigma_{1, j}^2), \, \mathcal{N}(\mu_{2, j}, \sigma_{2, j}^2) \right)
\end{equation}
\end{proof}

\subsection{One-Step Convergence of the Bures-Wasserstein Barycenter Iteration}
We now prove that when all expert covariances are diagonal, the general multi-dimensional fixed-point iteration for the Bures-Wasserstein barycenter reduces to a closed-form weighted sum that converges in exactly one iteration.

\begin{proposition}[One-Step Diagonal Convergence]
Let $\{\Sigma_k\}_{k=1}^K$ be a set of diagonal covariance matrices:
\begin{equation}
  \Sigma_k = \mathrm{diag}(\sigma_{k, 1}^2, \dots, \sigma_{k, d}^2) \quad \forall k
\end{equation}
Let $\Sigma_c = \mathrm{diag}(\sigma_{c, 1}^2, \dots, \sigma_{c, d}^2)$ be the diagonal Bures-Wasserstein barycenter. The fixed-point iteration is:
\begin{equation}
  \Sigma_c^{(t+1)} = \Sigma_c^{(t)-1/2} \left( \frac{1}{K} \sum_{k=1}^K \left( \Sigma_c^{(t)1/2} \Sigma_k \Sigma_c^{(t)1/2} \right)^{1/2} \right)^2 \Sigma_c^{(t)-1/2}
\end{equation}
If initialized with any positive diagonal covariance matrix $\Sigma_c^{(0)} = \mathrm{diag}(\sigma_{c, 1}^{(0)2}, \dots, \sigma_{c, d}^{(0)2})$ where $\sigma_{c, j}^{(0)} > 0$, then the iteration converges to the exact Bures-Wasserstein barycenter $\Sigma_c$ in exactly one iteration ($t = 0$), where:
\begin{equation}
  \sigma_{c, j} = \frac{1}{K} \sum_{k=1}^K \sigma_{k, j} \quad \forall j
\end{equation}
\end{proposition}

\begin{proof}
Let all matrices $\Sigma_c^{(t)}$ and $\Sigma_k$ be diagonal. The matrix operations commute, and the update equation simplifies coordinate-wise for each diagonal index $j \in \{1, \dots, d\}$:
\begin{equation}
  \sigma_{c, j}^{(t+1)2} = \frac{1}{\sigma_{c, j}^{(t)}} \left( \frac{1}{K} \sum_{k=1}^K \left( \sigma_{c, j}^{(t)} \sigma_{k, j}^2 \sigma_{c, j}^{(t)} \right)^{1/4} \right)^2 \frac{1}{\sigma_{c, j}^{(t)}}
\end{equation}
Since $\sigma_{c, j}^{(t)} > 0$ and $\sigma_{k, j} \ge 0$, we have:
\begin{equation}
  \left( \sigma_{c, j}^{(t)2} \sigma_{k, j}^2 \right)^{1/4} = \sqrt{\sigma_{c, j}^{(t)} \sigma_{k, j}}
\end{equation}
Substituting this back into the coordinate-wise iteration:
\begin{equation}
\begin{split}
  \sigma_{c, j}^{(t+1)2} &= \frac{1}{\sigma_{c, j}^{(t)2}} \left( \frac{1}{K} \sum_{k=1}^K \sqrt{\sigma_{c, j}^{(t)} \sigma_{k, j}} \right)^2 \\
  &= \frac{1}{\sigma_{c, j}^{(t)2}} \left( \sqrt{\sigma_{c, j}^{(t)}} \cdot \frac{1}{K} \sum_{k=1}^K \sigma_{k, j} \right)^2 \\
  &= \frac{1}{\sigma_{c, j}^{(t)2}} \cdot \sigma_{c, j}^{(t)} \cdot \left( \frac{1}{K} \sum_{k=1}^K \sigma_{k, j} \right)^2 \\
  &= \left( \frac{1}{K} \sum_{k=1}^K \sigma_{k, j} \right)^2
\end{split}
\end{equation}
This shows that for any $t \ge 0$, the updated standard deviation $\sigma_{c, j}^{(t+1)}$ is independent of the previous value $\sigma_{c, j}^{(t)}$:
\begin{equation}
  \sigma_{c, j}^{(t+1)} = \frac{1}{K} \sum_{k=1}^K \sigma_{k, j}
\end{equation}
Thus, starting from any positive diagonal initialization $\Sigma_c^{(0)}$, the first update $t = 0$ yields:
\begin{equation}
  \sigma_{c, j}^{(1)} = \frac{1}{K} \sum_{k=1}^K \sigma_{k, j} = \sigma_{c, j} \quad \forall j
\end{equation}
and subsequent updates remain stationary: $\Sigma_c^{(t)} = \Sigma_c^{(1)}$ for all $t \ge 1$. This proves that the iteration converges in exactly one step.
\end{proof}

\subsection{Frobenius Norm Contraction and Lipschitz Stability Guarantees}
We derive a formal contractive guarantee for the parameter magnitudes under Diagonal Bures-Wasserstein Parameter Alignment (D-BWPA), proving that our alignment operator does not amplify weight norms or destabilize the Lipschitz constant of the network.

\begin{proposition}[Frobenius Norm Contraction]
Let $W_{\mathrm{aligned}} \in \mathbb{R}^{d_{\mathrm{out}} \times d_{\mathrm{in}}}$ represent the aligned weight matrix under Diagonal Bures-Wasserstein Parameter Alignment (D-BWPA) for any regularization parameter $\epsilon \ge 0$. Let $\{W_{e, k}\}_{k=1}^K$ represent the weight matrices of the $K$ expert models. Then, the Frobenius norm of the aligned weights is strictly bounded by the average Frobenius norm of the expert weights:
\begin{equation}
  \|W_{\mathrm{aligned}}\|_F^2 \le \frac{1}{K} \sum_{k=1}^K \|W_{e, k}\|_F^2
\end{equation}
Furthermore, if $\epsilon > 0$ and the merged standard deviations $\sigma_{\mathrm{merged}, i} > 0$ for all $i$, the contraction is strict:
\begin{equation}
  \|W_{\mathrm{aligned}}\|_F^2 < \frac{1}{K} \sum_{k=1}^K \|W_{e, k}\|_F^2
\end{equation}
\end{proposition}

\begin{proof}
By definition, the squared Frobenius norm of a matrix $W \in \mathbb{R}^{d_{\mathrm{out}} \times d_{\mathrm{in}}}$ is the sum of the squared $L_2$ norms of its rows:
\begin{equation}
  \|W\|_F^2 = \sum_{i=1}^{d_{\mathrm{out}}} \|W_i\|^2
\end{equation}
For any expert $k$ and row $i$, the mean $\mu_{e, k, i} = \frac{1}{d_{\mathrm{in}}} \sum_{j=1}^{d_{\mathrm{in}}} W_{e, k, i, j}$ and standard deviation $\sigma_{e, k, i} = \sqrt{\frac{1}{d_{\mathrm{in}}} \sum_{j=1}^{d_{\mathrm{in}}} (W_{e, k, i, j} - \mu_{e, k, i})^2}$ define the row elements. Thus, the squared $L_2$ norm of the $i$-th row of the $k$-th expert's weight matrix is:
\begin{equation}
\begin{split}
  \|W_{e, k, i}\|^2 &= \sum_{j=1}^{d_{\mathrm{in}}} W_{e, k, i, j}^2 \\
  &= \sum_{j=1}^{d_{\mathrm{in}}} \left( \mu_{e, k, i} + (W_{e, k, i, j} - \mu_{e, k, i}) \right)^2 \\
  &= \sum_{j=1}^{d_{\mathrm{in}}} \mu_{e, k, i}^2 + \sum_{j=1}^{d_{\mathrm{in}}} (W_{e, k, i, j} - \mu_{e, k, i})^2 \\
  &\quad + 2 \mu_{e, k, i} \sum_{j=1}^{d_{\mathrm{in}}} (W_{e, k, i, j} - \mu_{e, k, i}) \\
  &= d_{\mathrm{in}} \mu_{e, k, i}^2 + d_{\mathrm{in}} \sigma_{e, k, i}^2 + 0 \\
  &= d_{\mathrm{in}} \left( \mu_{e, k, i}^2 + \sigma_{e, k, i}^2 \right)
\end{split}
\end{equation}
Summing this across all $d_{\mathrm{out}}$ rows yields the total Frobenius norm of the $k$-th expert matrix:
\begin{equation}
  \|W_{e, k}\|_F^2 = d_{\mathrm{in}} \sum_{i=1}^{d_{\mathrm{out}}} \left( \mu_{e, k, i}^2 + \sigma_{e, k, i}^2 \right)
\end{equation}

Now we analyze the aligned weights $W_{\mathrm{aligned}}$ under D-BWPA for any $\epsilon \ge 0$:
\begin{equation}
  W_{\mathrm{aligned}, i, j} = \mu_{\mathrm{target}, i} + \frac{\sigma_{\mathrm{target}, i}}{\sigma_{\mathrm{merged}, i} + \epsilon} \tilde{W}_{\mathrm{WA}, i, j}
\end{equation}
where $\tilde{W}_{\mathrm{WA}, i, j} = W_{\mathrm{WA}, i, j} - \mu_{\mathrm{merged}, i}$ represents the zero-mean components of the merged weights. By construction, we have:
\begin{equation}
  \sum_{j=1}^{d_{\mathrm{in}}} \tilde{W}_{\mathrm{WA}, i, j} = 0, \quad \sum_{j=1}^{d_{\mathrm{in}}} \tilde{W}_{\mathrm{WA}, i, j}^2 = d_{\mathrm{in}} \sigma_{\mathrm{merged}, i}^2
\end{equation}
The squared $L_2$ norm of the $i$-th row of the aligned weights is:
\begin{equation}
\begin{split}
  \|W_{\mathrm{aligned}, i}\|^2 &= \sum_{j=1}^{d_{\mathrm{in}}} W_{\mathrm{aligned}, i, j}^2 \\
  &= \sum_{j=1}^{d_{\mathrm{in}}} \left( \mu_{\mathrm{target}, i} + \frac{\sigma_{\mathrm{target}, i}}{\sigma_{\mathrm{merged}, i} + \epsilon} \tilde{W}_{\mathrm{WA}, i, j} \right)^2 \\
  &= \sum_{j=1}^{d_{\mathrm{in}}} \mu_{\mathrm{target}, i}^2 + \left( \frac{\sigma_{\mathrm{target}, i}}{\sigma_{\mathrm{merged}, i} + \epsilon} \right)^2 \sum_{j=1}^{d_{\mathrm{in}}} \tilde{W}_{\mathrm{WA}, i, j}^2 \\
  &\quad + 2 \mu_{\mathrm{target}, i} \frac{\sigma_{\mathrm{target}, i}}{\sigma_{\mathrm{merged}, i} + \epsilon} \sum_{j=1}^{d_{\mathrm{in}}} \tilde{W}_{\mathrm{WA}, i, j} \\
  &= d_{\mathrm{in}} \mu_{\mathrm{target}, i}^2 + \left( \frac{\sigma_{\mathrm{target}, i}}{\sigma_{\mathrm{merged}, i} + \epsilon} \right)^2 \left( d_{\mathrm{in}} \sigma_{\mathrm{merged}, i}^2 \right) + 0 \\
  &= d_{\mathrm{in}} \left( \mu_{\mathrm{target}, i}^2 + \sigma_{\mathrm{target}, i}^2 \frac{\sigma_{\mathrm{merged}, i}^2}{(\sigma_{\mathrm{merged}, i} + \epsilon)^2} \right)
\end{split}
\end{equation}
Observe that for any $\epsilon \ge 0$, we have $\frac{\sigma_{\mathrm{merged}, i}^2}{(\sigma_{\mathrm{merged}, i} + \epsilon)^2} \le 1$. Therefore:
\begin{equation}
  \|W_{\mathrm{aligned}, i}\|^2 \le d_{\mathrm{in}} \left( \mu_{\mathrm{target}, i}^2 + \sigma_{\mathrm{target}, i}^2 \right)
\end{equation}
where the inequality is strict ($<$) whenever $\epsilon > 0$ and $\sigma_{\mathrm{merged}, i} > 0$.
Summing across all $d_{\mathrm{out}}$ rows yields the total Frobenius norm of the aligned weights:
\begin{equation}
  \|W_{\mathrm{aligned}}\|_F^2 \le d_{\mathrm{in}} \sum_{i=1}^{d_{\mathrm{out}}} \left( \mu_{\mathrm{target}, i}^2 + \sigma_{\mathrm{target}, i}^2 \right)
\end{equation}

Recall that the target statistics in D-BWPA are defined as the averages of the expert statistics:
\begin{equation}
  \mu_{\mathrm{target}, i} = \frac{1}{K} \sum_{k=1}^K \mu_{e, k, i}, \quad \sigma_{\mathrm{target}, i} = \frac{1}{K} \sum_{k=1}^K \sigma_{e, k, i}
\end{equation}
Since the function $g(x) = x^2$ is convex on $\mathbb{R}$, we apply Jensen's inequality to obtain:
\begin{align}
  \mu_{\mathrm{target}, i}^2 &= \left( \frac{1}{K} \sum_{k=1}^K \mu_{e, k, i} \right)^2 \le \frac{1}{K} \sum_{k=1}^K \mu_{e, k, i}^2 \\
  \sigma_{\mathrm{target}, i}^2 &= \left( \frac{1}{K} \sum_{k=1}^K \sigma_{e, k, i} \right)^2 \le \frac{1}{K} \sum_{k=1}^K \sigma_{e, k, i}^2
\end{align}
Adding these two inequalities together:
\begin{equation}
  \mu_{\mathrm{target}, i}^2 + \sigma_{\mathrm{target}, i}^2 \le \frac{1}{K} \sum_{k=1}^K \left( \mu_{e, k, i}^2 + \sigma_{e, k, i}^2 \right)
\end{equation}
Summing this row-wise bound across all $d_{\mathrm{out}}$ rows and multiplying by $d_{\mathrm{in}}$:
\begin{equation}
\begin{split}
  \|W_{\mathrm{aligned}}\|_F^2 &= d_{\mathrm{in}} \sum_{i=1}^{d_{\mathrm{out}}} \left( \mu_{\mathrm{target}, i}^2 + \sigma_{\mathrm{target}, i}^2 \right) \\
  &\le d_{\mathrm{in}} \sum_{i=1}^{d_{\mathrm{out}}} \left( \frac{1}{K} \sum_{k=1}^K \left( \mu_{e, k, i}^2 + \sigma_{e, k, i}^2 \right) \right) \\
  &= \frac{1}{K} \sum_{k=1}^K \left( d_{\mathrm{in}} \sum_{i=1}^{d_{\mathrm{out}}} \left( \mu_{e, k, i}^2 + \sigma_{e, k, i}^2 \right) \right) \\
  &= \frac{1}{K} \sum_{k=1}^K \|W_{e, k}\|_F^2
\end{split}
\end{equation}
which completes the proof.
\end{proof}

\subsection{Rademacher Complexity and Generalization Guarantees under D-BWPA}
Using our Frobenius norm contraction guarantee, we establish formal generalization bounds for the merged model under Diagonal Bures-Wasserstein Parameter Alignment (D-BWPA). We leverage standard spectral-normalized Rademacher complexity bounds for deep neural networks \cite{bartlett2017spectrally} to demonstrate that merging via D-BWPA mathematically guarantees stability against overfitting and maintains rigorous generalization capacity.

\begin{proposition}[Rademacher Complexity Bounds]
Let $\mathcal{H}_{d_{\mathrm{out}}, d_{\mathrm{in}}}^{L}$ represent the class of $L$-layer feedforward neural networks mapping input $x \in \mathcal{X} \subset \mathbb{R}^{d_0}$ (with $\|x\|_2 \le B_x$) to output $y \in \mathbb{R}^{d_L}$. Let $W^{(l)} \in \mathbb{R}^{d_l \times d_{l-1}}$ represent the weight matrix at layer $l \in \{1, \dots, L\}$. Under D-BWPA, the merged network's Rademacher complexity $\mathcal{R}_S(\mathcal{H}_{\mathrm{aligned}})$ on a sample $S$ of size $m$ is bounded by:
\begin{equation}
  \mathcal{R}_S(\mathcal{H}_{\mathrm{aligned}}) \le \frac{2 B_x \sqrt{L + 1} \ln(2d_{\max}) \prod_{l=1}^L \left( \frac{1}{K} \sum_{k=1}^K \|W_{e, k}^{(l)}\|_F^2 \right)^{1/2}}{\sqrt{m}}
\end{equation}
where $W_{e, k}^{(l)}$ is the weight matrix of expert $k$ at layer $l$, and $d_{\max} = \max_l d_l$ represents the maximum layer dimension.
\end{proposition}

\begin{proof}
According to the spectral-normalized Rademacher complexity bound for deep neural networks with Lipschitz-continuous activations (such as ReLU) of \cite{bartlett2017spectrally}, the empirical Rademacher complexity of a network class with layer-wise Frobenius norm bounds $\|W^{(l)}\|_F \le B^{(l)}$ is bounded by:
\begin{equation}
  \mathcal{R}_S(\mathcal{H}) \le \frac{2 B_x \sqrt{L + 1} \ln(2d_{\max}) \prod_{l=1}^L B^{(l)}}{\sqrt{m}}
\end{equation}
Under D-BWPA (for any $\epsilon \ge 0$), we proved in Proposition A.3 (Frobenius Norm Contraction) that for any layer $l \in \{1, \dots, L\}$:
\begin{equation}
  \|W_{\mathrm{aligned}}^{(l)}\|_F \le \left( \frac{1}{K} \sum_{k=1}^K \|W_{e, k}^{(l)}\|_F^2 \right)^{1/2}
\end{equation}
Substituting this contractive bound as the layer-wise Frobenius norm bound $B_{\mathrm{aligned}}^{(l)} = \left( \frac{1}{K} \sum_{k=1}^K \|W_{e, k}^{(l)}\|_F^2 \right)^{1/2}$ for each layer $l$ of the aligned merged model directly yields:
\begin{equation}
  \mathcal{R}_S(\mathcal{H}_{\mathrm{aligned}}) \le \frac{2 B_x \sqrt{L + 1} \ln(2d_{\max}) \prod_{l=1}^L \left( \frac{1}{K} \sum_{k=1}^K \|W_{e, k}^{(l)}\|_F^2 \right)^{1/2}}{\sqrt{m}}
\end{equation}
This bound guarantees that the Rademacher complexity of the D-BWPA merged model is strictly controlled by the root-mean-squared Frobenius norms of the individual expert models. This mathematically guarantees that if the expert models are well-regularized (i.e., have bounded weight norms), the merged model will also preserve generalization capabilities and be protected against catastrophic representation explosion or variance expansion.
\end{proof}

\section{Experimental Settings and Reproducibility Details}
\label{sec:appendix_experimental}

To ensure complete reproducibility of our empirical results, we document the precise model architectures, preprocessing pipelines, training hyperparameters, and calibration procedures.

\subsection{Model Architectures}
We evaluated our methods on two distinct deep neural network families:
\begin{enumerate}
  \item \textbf{ResNet-18 Expert:} We adopted PyTorch's standard 18-layer Residual Network (ResNet-18) backbone. The fully connected classification layer was replaced with an identity mapping, and a new linear head with $10$ outputs was added:
  \begin{equation}
    f(x) = \mathbf{W}_{\mathrm{head}} \cdot \mathrm{Backbone}(x)
  \end{equation}
  where the backbone outputs a $512$-dimensional representation. The total number of parameters in the ResNet-18 expert is approximately $11.18$ million.
  \item \textbf{MLP Expert (No BatchNorm):} The Multi-Layer Perceptron (MLP) expert comprises a 3-layer feedforward neural network structure:
  \begin{itemize}
    \item Input Layer: Linear transformation from $3072$ dimensions (for flattened $32 \times 32 \times 3$ images) to $512$ channels, followed by a ReLU activation.
    \item Hidden Layer 1: Linear transformation from $512$ to $512$ channels, followed by a ReLU activation.
    \item Hidden Layer 2: Linear transformation from $512$ to $512$ channels, followed by a ReLU activation.
    \item Output Head: Linear classification layer mapping the $512$-dimensional representation to $10$ logits.
  \end{itemize}
  The total number of parameters in the MLP expert is approximately $2.11$ million. This network strictly contains no BatchNorm or other normalization layers to evaluate merging stability under pure feedforward propagation.
\end{enumerate}

\subsection{Training and Fine-Tuning Hyperparameters}
All task-specific experts were initialized from a shared progenitor model (pretrained ResNet-18 weights for ResNet-18, and a randomly initialized common seed for MLP).
For each of the three tasks—MNIST, FashionMNIST, and CIFAR-10—we fine-tuned the expert models independently using the hyperparameters detailed in Table~\ref{tab:hyperparameters}.

\begin{table}[h]
\caption{Expert Training Hyperparameters.}
\label{tab:hyperparameters}
\vskip 0.1in
\begin{center}
\begin{small}
\begin{sc}
\begin{tabular}{lc}
\toprule
Hyperparameter & Value \\
\midrule
Optimizer & AdamW \\
Learning Rate & $5 \times 10^{-4}$ \\
Weight Decay & $1 \times 10^{-2}$ \\
Batch Size & $128$ \\
Fine-Tuning Epochs & $5$ \\
Hardware & 4 CPUs, NVIDIA H100 GPU \\
Loss Function & Cross-Entropy \\
\bottomrule
\end{tabular}
\end{sc}
\end{small}
\end{center}
\end{table}

\subsection{Data Preprocessing and Augmentation}
All images are preprocessed and scaled to ensure uniform tensor dimensions across MNIST, FashionMNIST, and CIFAR-10:
\begin{itemize}
  \item Image Resize: All input images are bilinearly interpolated to size $32 \times 32$ pixels.
  \item Channel Alignment: Single-channel greyscale images (MNIST and FashionMNIST) are repeated across 3 channels to match the CIFAR-10 format.
  \item Normalization: Tensors are normalized using channel-wise means of $[0.5, 0.5, 0.5]$ and standard deviations of $[0.5, 0.5, 0.5]$, scaling input values to $[-1.0, 1.0]$.
\end{itemize}

\subsection{BatchNorm Calibration Details}
For ResNet-18, Data-Efficient BatchNorm Calibration (DE-BN) was executed using a single batch of $32$ calibration samples ($11$ or $10$ samples drawn randomly from the training split of each of the three tasks to prevent distribution bias). Our momentum-corrected calibration algorithm (Algorithm~\ref{alg:debn}) forces the BN momentum to $\beta = 1.0$ during the calibration forward pass, ensuring the running statistics precisely equal the batch statistics before reverting to the default $\beta = 0.1$ for evaluation.

\end{document}
